% =============================================================================
\section{Network Graph Representation}
\label{sec:network}
% =============================================================================

The multi-tank system is described as a directed graph
\begin{equation}
  G = (\Nodes,\, \Edges),
\end{equation}
where $\Nodes$ is the set of \emph{nodes} and $\Edges \subseteq \Nodes\times\Nodes$
is the set of directed \emph{edges}.  Every mass-transfer path in the system is
represented by an edge; the governing equations (Section~\ref{sec:dae}) are then
assembled purely from graph incidence.

% -----------------------------------------------------------------------------
\subsection{Nodes}
\label{sec:network:nodes}
% -----------------------------------------------------------------------------

Three node types are recognised:

\begin{center}
\begin{tabularx}{\linewidth}{@{} l l X @{}}
\toprule
\textbf{Type} & \textbf{YAML \texttt{type}} & \textbf{Description} \\
\midrule
\texttt{tank}       & \texttt{tank}       &
  Finite-volume hydrogen storage vessel with evolving thermodynamic state
  $(m_i, T_i, T_{s,i})$. \\[4pt]
\texttt{source}     & \texttt{source}     &
  Infinite reservoir that supplies mass at a prescribed enthalpy (e.g.\ a
  cryogenic pump feeding a refuel edge). \\[4pt]
\texttt{sink}       & \texttt{sink}       &
  Infinite reservoir that absorbs mass with no back-pressure (e.g.\
  \texttt{atmosphere} for venting, \texttt{fuel\_cell} for mission discharge). \\
\bottomrule
\end{tabularx}
\end{center}

Only \texttt{tank} nodes carry dynamic states; \texttt{source} and \texttt{sink}
nodes act as boundary conditions.

% -----------------------------------------------------------------------------
\subsection{Edges}
\label{sec:network:edges}
% -----------------------------------------------------------------------------

Four edge types are currently implemented, summarised in
\cref{tab:edge-types}.

\begin{table}[ht]
\centering
\caption{Edge types, their activation conditions, and the enthalpy assigned to
the transferred mass.}
\label{tab:edge-types}
\begin{tabularx}{\linewidth}{@{} l l X l @{}}
\toprule
\textbf{Type} & \textbf{Direction} & \textbf{Activation / flow rate} & \textbf{Enthalpy} \\
\midrule
\texttt{coupling}  & tank $\to$ tank &
  Activated by a valve model (e.g.\ pressure-triggered, PID, governor)
  when the pressure differential across the edge meets the opening criterion. &
  $h_\text{source}(T_\text{src},\rho_\text{src})$ \\[6pt]

\texttt{discharge} & tank $\to$ sink &
  Flow rate prescribed by a mission CSV file, $\mdot_\text{CSV}(t)$.
  In Config~B, throttled to the algebraic limit (see \cref{sec:dae:configB}). &
  $h(T_i,\rho_i)$ \\[6pt]

\texttt{vent}      & tank $\to$ sink &
  Activated when $P_i \ge P_{\text{vent},i}$.
  Proportional rate model: $\mdot_\text{vent} = k_v\max(0, P_i - P_{\text{vent},i})$. &
  $h(T_i,\rho_i)$ \\[6pt]

\texttt{refuel}    & source $\to$ tank &
  Flow rate prescribed by a mission CSV or fixed schedule.
  Enthalpy computed from cryogenic pump isentropic model. &
  $h_\text{cryo}(P_i)$ \\
\bottomrule
\end{tabularx}
\end{table}

% -----------------------------------------------------------------------------
\subsection{YAML Schema Mapping}
\label{sec:network:yaml}
% -----------------------------------------------------------------------------

The graph is declared in the \texttt{network} block of the scenario YAML file:

\begin{lstlisting}[language=Python, caption={Excerpt of YAML network declaration.}]
network:
  nodes:
    - node_id: 1
      type: tank          # tank / source / sink
      fluid: CH2
      ...
    - node_id: atmosphere
      type: sink

  edges:
    - edge_id: ch2_cch2_coupling
      connection_type: coupling     # coupling / discharge / vent / refuel
      from_node: 1
      to_node: 2
      ...
    - edge_id: discharge_fc
      connection_type: discharge
      from_node: 2
      to_node: fuel_cell
      ...
\end{lstlisting}

Mission discharge and venting edges that were previously implicit (hard-coded in
the dynamic model) are now explicit edges in the \texttt{network.edges} list.

% -----------------------------------------------------------------------------
\subsection{Example: CH2/CCH2 Two-Tank System}
\label{sec:network:example}
% -----------------------------------------------------------------------------

\Cref{fig:network-example} shows the directed graph for the reference scenario:
a 800-bar compressed hydrogen tank (CH2) coupled to a 400-bar cryo-compressed
hydrogen tank (CCH2) with mission discharge assigned to Tank~2.

\begin{figure}[ht]
\centering
\begin{tikzpicture}[
    node distance = 2.2cm and 3.0cm,
    tank/.style    = {draw, rounded corners=4pt, fill=blue!8, minimum width=2.8cm,
                      minimum height=1.2cm, align=center, font=\small\bfseries},
    boundary/.style= {draw, dashed, rounded corners=4pt, fill=gray!10,
                      minimum width=2.0cm, minimum height=0.9cm, align=center,
                      font=\small\itshape},
    edge/.style    = {-{Latex[length=6pt]}, thick},
    label/.style   = {font=\scriptsize, midway, align=center, text width=2.2cm},
  ]

  % Nodes
  \node[tank]     (t1)                      {Tank 1\\CH2\\800 bar};
  \node[tank]     (t2) [right=4.5cm of t1]  {Tank 2\\CCH2\\400 bar};
  \node[boundary] (fc) [right=3.5cm of t2]  {Fuel Cell\\(sink)};
  \node[boundary] (atm1) [above=1.8cm of t1]{Atmosphere\\(sink)};
  \node[boundary] (atm2) [above=1.8cm of t2]{Atmosphere\\(sink)};

  % Coupling edge
  \draw[edge, blue!70!black]
        (t1) -- node[label, above] {coupling\\$P_1 > P_2 + \Delta P$} (t2);

  % Vent edges
  \draw[edge, red!70!black]
        (t1) -- node[label, right] {vent\\$P_1 \ge P_{\text{vent},1}$} (atm1);
  \draw[edge, red!70!black]
        (t2) -- node[label, right] {vent\\$P_2 \ge P_{\text{vent},2}$} (atm2);

  % Discharge edge
  \draw[edge, green!50!black]
        (t2) -- node[label, above] {discharge\\$\mdot_\text{CSV}(t)$} (fc);

\end{tikzpicture}
\caption{Directed network graph for the CH2/CCH2 coupled system.  Tank~1 has no
discharge edge; Tank~2 carries the full mission discharge load.  Both tanks have
vent edges to the atmosphere sink.}
\label{fig:network-example}
\end{figure}
